\documentclass[11pt]{article}
\usepackage{amsmath,amssymb,amsthm}
\usepackage[margin=1in]{geometry}
\usepackage[colorlinks=true,linkcolor=blue,citecolor=blue,urlcolor=blue]{hyperref}

\newtheorem{theorem}{Theorem}
\DeclareMathOperator{\Ric}{Ric}
\DeclareMathOperator{\Hess}{Hess}
\DeclareMathOperator{\grad}{grad}
\newcommand{\dS}{\mathrm{dS}}
\newcommand{\Lie}{\mathcal L}
\title{Nariai rigidity for null-complete vacuum spacetimes\\with a semidefinite Cauchy slice}
\author{}
\date{}

\begin{document}
\maketitle

\begin{abstract}
Let a four-dimensional Einstein spacetime with positive cosmological
constant be globally hyperbolic and null geodesically complete in both
directions. Suppose that it has a compact nonspherical Cauchy hypersurface
with positive semidefinite second fundamental form. The universal
Lorentzian cover is the product of the universal cover of two-dimensional
de Sitter space and a round two-sphere, both of curvature equal to the
cosmological constant. The proof uses the two totally geodesic null
foliations of Galloway and Ling, but does not infer a metric product from
the foliations alone. A stable compact initial cut makes their horizontal
distribution integrable nearby. The vacuum equations and smoothness at
the extrema of a putative warp factor force that factor to be constant.
Local product geometry then propagates through noncompact as well as
compact foliation charts.
\end{abstract}

\section{Statement and geometric input}

We use Lorentzian signature $(-+++)$ and take the second fundamental form
of a spacelike Cauchy hypersurface with future unit normal $T$ to be
$K(X,Y)=g(\nabla_X T,Y)$. A spherical space means a three-manifold
diffeomorphic to a quotient of the round three-sphere by a finite group
acting freely by isometries.

\begin{theorem}\label{thm:main}
Let $(M,g)$ be a connected smooth time-oriented globally hyperbolic
four-dimensional spacetime with $\Ric_g=\Lambda g$, $\Lambda>0$.
Suppose every null geodesic is complete in both affine directions.
If $M$ has a smooth compact spacelike Cauchy hypersurface $V$ with
$K\geq 0$ and $V$ is not a spherical space, then its universal
Lorentzian cover is isometric to
\[
   \widetilde{\dS_2(\Lambda)}\times S^2(\Lambda).
\]
Here each factor has constant sectional curvature $\Lambda$. The
statement with $K\leq 0$ follows by reversing the time orientation.
\end{theorem}

A classification of the deck transformations is not part of the theorem.
The topological reduction and existence of the two null foliations come
from Galloway--Ling \cite{GLcos,GLpen}. We specify the initial section of
these foliations because that section supplies essential metric
information. In particular, the argument does not require saturation of
a stable marginally outer trapped surface area bound, an assumption
used in the separate rigidity setting of \cite{GM}.

By the topological reduction preceding \cite[Theorem~5]{GLcos},
a finite cover has an orientable compact Cauchy hypersurface with
nontrivial second homology. Work first in that cover. Minimize area in
a nonzero integral second homology class and choose a connected
component carrying a nonzero class. In dimension three this gives a
smooth closed oriented two-sided stable minimal surface $\Sigma$.
It is nonseparating, since a separating oriented component would
represent the zero class. The signs of both future null expansions
are nonnegative: the spatial mean curvature is zero and
$\operatorname{tr}_{T\Sigma}K\geq0$. The null energy condition used
in \cite{GLcos,GLpen} holds since $\Ric_g(k,k)=0$ for null $k$.

We now use the cyclic-cover construction in the proofs of
\cite[Theorems~8--9]{GLpen}, starting from this particular $\Sigma$.
In the spacetime cover there are transverse foliations by totally
geodesic null hypersurfaces. The two leaves through the chosen
lift of $\Sigma$ have that lift as their intersection locally along
it. To see which cut is at issue, the two families on the lifted
Cauchy slice start from the deck translates of the chosen surface;
the proof of Theorem~9 first rules out a strict positive null
expansion on that surface and then takes the two null normal
congruences of the resulting marginal families. We use the smooth
leaf charts in the cover, not achronality of the projected leaves
in the original spacetime. For clarity, smoothness transverse to a
complete generator can also be checked on the geodesic map from an
initial family of cuts. Its screen differential retains rank because
the null second fundamental form vanishes and the screen metric is
transported unchanged. A Jacobi variation $J$ of the initial cut
parameter satisfies
$k[g(J,k)]=g(\nabla_k J,k)=\tfrac12 J[g(k,k)]=0$
for the affinely parameterized generator $k$. This scalar is nonzero
initially, by transversality in the Cauchy slice, so the remaining
transverse differential retains rank. All arguments below are on the
connected cyclic spacetime cover unless otherwise stated. All covers
here inherit null completeness by geodesic path lifting.

\section{The stable cut and the vacuum curvature identities}

Write $\nu$ for the unit normal of the lifted $\Sigma$ in the Cauchy
slice. Both null second fundamental forms vanish on $\Sigma$.
Their sum and difference give $K|_{T\Sigma\times T\Sigma}=0$ and
$\mathrm{II}^{V}_\Sigma=0$. Positive semidefiniteness then gives
$K(\nu,X)=0$ for $X\in T\Sigma$: a vector $X$ on which the quadratic
form vanishes is in its kernel. Thus $K$ has rank at most one there.
The Hamiltonian constraint reduces on $\Sigma$ to $R_V=2\Lambda$.
The constant-function stability test gives
$\int_\Sigma \Ric_V(\nu,\nu)\,dA\leq0$, since
$\mathrm{II}^{V}_\Sigma=0$. The Gauss equation yields
\[
 \int_\Sigma R_\Sigma\,dA
   =\int_\Sigma (2\Lambda-2\Ric_V(\nu,\nu))\,dA
   \ \geq\ 2\Lambda\operatorname{Area}(\Sigma)>0.
\]
Consequently $\Sigma$ is a two-sphere by Gauss--Bonnet. In the normalized
null frame
\[
 \ell=(T+\nu)/\sqrt2,\qquad n=(T-\nu)/\sqrt2,\qquad g(\ell,n)=-1 ,
\]
the normal connection $\zeta(X)=g(\nabla_X\ell,n)$ is
$-K(X,\nu)=0$ on $\Sigma$. Although a boost of the null frame can
change $\zeta$ by an exact form, the curvature $d\zeta$ is independent
of this choice.

Here are the local vacuum identities we need. In a double-null chart,
use an adapted null tetrad $(\ell,n,m,\bar m)$; the complex dyad is
tangent to the spacelike plaques cut out by the two null foliations.
Both null congruences are geodesic, with zero expansion, shear and
optical twist. The four Weyl scalars
$\Psi_0,\Psi_1,\Psi_3,\Psi_4$ vanish. Indeed curvature preserves a
totally geodesic null tangent hyperplane when all three input vectors
are tangent. Contraction with its null normal $\ell$ shows that the
screen components $C_{\ell A\ell B}$ and $C_{\ell ABC}$ vanish,
where $C$ denotes the Weyl tensor; the constant-curvature part of
the Einstein Riemann tensor vanishes in these components. By
trace-freeness $C_{\ell n\ell A}$ also vanishes. Repeat for $n$.
The surviving scalar $\Psi_2$ is invariant under the adapted boosts
and dyad rotations. Reversal of screen orientation conjugates it.

We use the Newman--Penrose Bianchi identities in the conventions
of \cite[Appendix A, Eqs.~(A24), (A25), (A28), (A29)]{NP}. In the
present Einstein case the Ricci derivative terms vanish and the
above optical conditions give
\begin{equation}\label{eq:bianchi}
 D\Psi_2=0,\qquad \Delta_{\rm NP}\Psi_2=0,\qquad
 \delta\Psi_2=3\tau\Psi_2,\qquad
 \bar\delta\Psi_2=-3\pi\Psi_2 .
\end{equation}
Here $D,\Delta_{\rm NP},\delta,\bar\delta$ differentiate along the
four tetrad vectors. The coefficients $\tau,\pi$ are, up to the
tetrad sign conventions, the complex screen components of
$\nabla_n\ell$ and $\nabla_\ell n$, respectively. In particular,
each vanishes precisely when its indicated real screen projection
vanishes. The identities do not involve division by $\Psi_2$.
The normal Ricci equation, with both null second fundamental forms
zero, identifies $d\zeta$ with the mixed Riemann component
$R_{AB\ell n}$ up to sign. Its Einstein constant-curvature part is
zero and its Weyl part is proportional to $\operatorname{Im}\Psi_2$
times the screen area form, by the adapted tetrad decomposition.
In particular $\operatorname{Im}\Psi_2=0$ on $\Sigma$.

\section{Removal of the horizontal twist}

Consider first a sufficiently small double-null bundle box around the
whole compact $\Sigma$. The leaf labels give its base coordinates
$(u,v)$. The orthogonal complement of the spacelike plaques is
Lorentzian; let $U,V_0$ be the unique horizontal lifts of
$\partial_u,\partial_v$, respectively. These are null and we write
$g(U,V_0)=-\Omega<0$. The induced plaque metric is $h$.
The two null second fundamental forms vanish, so
$\Lie_Uh=\Lie_{V_0}h=0$ on vertical arguments. Both fields are
projectable, so $Z=[U,V_0]$ is vertical, and
$\Lie_Z h=[\Lie_U,\Lie_{V_0}]h=0$ on the same arguments. Thus $Z$
is Killing on each plaque.

For a vertical $X$, the Koszul formula gives
\[
 2g(\nabla_X U,V_0)=-X\Omega-h(Z,X).
\]
With $\ell=U/\sqrt{\Omega}$ and $n=V_0/\sqrt{\Omega}$ the derivative
of $\Omega$ cancels, so
\begin{equation}\label{eq:koszul}
 \zeta(X)=-\frac{h(Z,X)}{2\Omega},\qquad
 (\nabla_U V_0)^{\rm vert}=\tfrac12(\grad_h\Omega+Z),\qquad
 (\nabla_{V_0} U)^{\rm vert}=\tfrac12(\grad_h\Omega-Z).
\end{equation}
For example, in the middle identity the only terms in
$2g(\nabla_U V_0,X)$ that survive are $-Xg(U,V_0)$ and
$g(X,[U,V_0])$; brackets with a vertical field are vertical.

Here the compact box can be built by transporting $\Sigma$ with $U$
and then $V_0$ along small base segments. The two global leaf labels
on the starting sphere may be taken from the two initial families
in the cyclic construction. Compactness gives a uniform flow
interval around the whole sphere. For fixed labels the transport is a
diffeomorphism onto its image with inverse given by the reversed
segments. The resulting family is a neighborhood of $\Sigma$,
foliated there by compact spheres. By the first two equations
in \eqref{eq:bianchi}, $\operatorname{Im}\Psi_2$ vanishes throughout
this neighborhood. The normal Ricci equation therefore gives
$d\zeta=0$ on each sphere. As $H^1(S^2;\mathbb R)=0$, there is a
smooth potential $\varphi$ on each such cut with $\zeta=d\varphi$.
The divergence of its Killing field $Z$ is zero and the cut has no
boundary. Thus \eqref{eq:koszul} gives
\[
    0=\int Z(\varphi)\,dA_h
      =-\int\frac{|Z|_h^2}{2\Omega}\,dA_h .
\]
Hence $Z=0$ on the compact-cut box. This integral is the step
that removes horizontal twist. It uses initial normal-connection
curvature and compact cut topology, not simply the vanishing optical
twist of each individual congruence.

\section{The smooth compact warp is constant}

Now $U$ and $V_0$ commute, so horizontal flow trivializes the box
and $g=h(x)-2\Omega(u,v,x)\,du\,dv$, with $h$ independent of $(u,v)$.
Write $w=\tfrac12\log\Omega$. For base indices $a,b$ and plaque
indices $A,B$, the mixed Christoffel terms in the block metric are
\[
 \Gamma^a{}_{bA}=\delta^a_b\partial_A w,\qquad
 \Gamma^A{}_{ab}=-g_{ab}h^{AB}\partial_B w,\qquad
 \Gamma^A{}_{aB}=0=\Gamma^a{}_{AB}.
\]
As $\Gamma^b{}_{ab}=2\partial_a w$, the cross terms quadratic in
Christoffel symbols cancel in $\Ric_{aA}$ and the derivative terms
give $\Ric_{aA}=-\partial_a\partial_A w$. Thus in the Einstein
case $w=w_0(u,v)+w_1(x)$, by integrating the mixed identities on
the connected cut. Set $f=e^{w_1}>0$ and absorb $w_0$ into a
two-dimensional Lorentzian metric $g_B$. The full metric is
$h+f^2g_B$. The warped Ricci equations (obtained by substituting
the same block Christoffel terms) are
\begin{equation}\label{eq:einstein-warp}
 \Ric_h-2f^{-1}\Hess_h f=\Lambda h,\qquad
 \Ric_{g_B}=(\Lambda f^2+f\Delta_h f+|df|_h^2)g_B ,
\end{equation}
where $\Delta_h=\operatorname{div}_h\grad_h$. For example the
horizontal curvature contraction is that of $g_B$ minus
$(f\Delta_h f+|df|_h^2)g_B$; the mixed sectional terms give
$-2f^{-1}\Hess_h f$ on the cut since the other factor has dimension
two. In dimension two \eqref{eq:einstein-warp} implies
\begin{equation}\label{eq:warp-id}
 \Hess_h f=\tfrac12\Delta_h f\,h,\qquad
 K_h=\Lambda+\frac{\Delta_h f}{f},\qquad
 d|df|_h^2=\Delta_h f\,df,\qquad d\Delta_h f=-2K_h\,df .
\end{equation}
The last equality follows by taking the divergence of the Hessian
equation and commuting covariant derivatives:
$\tfrac12 d\Delta_h f=d\Delta_h f+K_h\,df$.

For the compact connected sphere, the curvature of $g_B$ is a
constant $\kappa$: the multiplier in the second equation of
\eqref{eq:einstein-warp} is a function of $x$ alone and the curvature
of $g_B$ a function of the base alone. By \eqref{eq:warp-id},
\[
 d(f|df|_h^2)=(|df|_h^2+f\Delta_h f)\,df
              =(\kappa-\Lambda f^2)\,df.
\]
Thus $|df|_h^2=Q(f)$, where $Q(s)=\kappa-\Lambda s^2/3+c/s$
and $c$ is a single constant on the cut. Substitution into
\eqref{eq:einstein-warp} gives $\Delta_h f=Q'(f)$ everywhere,
including at critical points.

If the minimum $a$ and maximum $b$ of $f$ satisfy $a<b$, they
are positive roots of $Q$. Neither $Q'(a)$ nor $Q'(b)$ can
vanish: along any unit cut geodesic $f$ satisfies
$(f\circ\gamma)''=Q'(f\circ\gamma)/2$ by the Hessian equation.
A critical point with $Q'=0$ gives equilibrium initial data, hence
constancy by uniqueness, first in a normal ball and then on the
connected cut. Therefore $Q'(a)>0>Q'(b)$.

Let $J$ be the parallel oriented rotation by $\pi/2$ on the
cut. The field $J\grad_h f$ is Killing, since its covariant
derivative is $(Q'(f)/2)J$. Its flow exists for all time and
fixes a minimum and a maximum. At these points the derivatives of
the flow are rotations with signed speeds $Q'(a)/2$ and
$Q'(b)/2$. An isometry of a connected Riemannian manifold fixing
a point and with identity derivative there is the identity:
it fixes a normal neighborhood, and the same argument continues
along paths. As a result, the smallest positive period of the
same global flow computed at either fixed point must be the
same. In particular $Q'(a)=-Q'(b)$. At any positive root $s$,
$Q'(s)=\kappa/s-\Lambda s$, so this equality requires
$\kappa=\Lambda ab$. On the other hand the two equations
$sQ(s)=0$ imply $\kappa=\Lambda(a^2+ab+b^2)/3$. They contradict
$a<b$ because $\Lambda>0$. Thus $f$ is constant. This period
argument enforces smoothness at both extrema without introducing
rotational coordinates or permitting conical defects.

By \eqref{eq:einstein-warp}, both factors, with the constant
$f^2$ absorbed into $g_B$, have curvature $\Lambda$. This is local
Nariai geometry throughout the initial neighborhood.

\section{Local continuation without compact late cuts}

Let $\mathcal P$ be the open set of points with a locally Nariai
neighborhood. It is nonempty, and we show it is closed. Denote by
$\Psi_*$ the real nonzero Nariai value of $\Psi_2$ in the conventions
above. Its nonvanishing follows, for example, from the spacelike
sectional curvature $\Lambda$ of the spherical factor compared
with the Einstein constant-curvature contribution $\Lambda/3$.
The adapted directions in local Nariai are those of the Lorentzian
factor: the Weyl tensor is nonzero of type D and these are precisely
its two repeated principal null directions.

Let $p\in\overline{\mathcal P}$. Choose a connected double-null
chart with a base rectangle and disk plaques around $p$. Shrink
the base rectangle so that transport along two base coordinate
segments from a smaller disk remains in the original chart; this
follows from the smooth bounded horizontal coordinate velocities
on a relatively compact box. Choose a small open plaque patch
in $\mathcal P$ at a point sufficiently close to $p$, within the
smaller disk. By transporting that patch along $U$ and then $V_0$
we obtain an open family meeting every plaque of the shrunk box.
Equation \eqref{eq:bianchi} preserves $\Psi_2=\Psi_*$ under
these transports. The two screen derivatives of $\Psi_2$ vanish
on this open family; the last two equations of \eqref{eq:bianchi}
give $\tau=\pi=0$ there. By \eqref{eq:koszul}, this gives both
$Z=0$ and $\grad_h\Omega=0$ on the transported patches.

On each connected plaque, $Z$ is Killing and is zero on a nonempty
open patch, so it vanishes on the entire plaque. For this uniqueness,
the value and first derivative of a Killing field determine it
along a geodesic by the Jacobi equation; continuation through
normal neighborhoods then applies even to an incomplete plaque.
Thus $Z=0$ in the box. Pass to commuting horizontal coordinates,
restricting to connected plaque disks containing $p$ (on its
plaque) and the transported patches. More explicitly, the initial
patch can be chosen arbitrarily close to $p$ in the original chart.
After shrinking the base, the horizontal trivialization and its
inverse have uniformly small displacements, and fixed inner
disk plaques can contain the smaller patch and the disk coordinate
of $p$. In these coordinates the mixed Einstein equation again
gives the warped form of \eqref{eq:einstein-warp}. On every
connected plaque $\grad_h f=0$ and $\Delta_h f=0$ on a nonempty
open patch because $\Omega$ has zero screen derivative there.
Along any smooth plaque path $\gamma$, equations
\eqref{eq:warp-id} are a first-order system for
$(f,df,\Delta_h f)$:
\[
 \frac{df}{ds}=df(\dot\gamma),\qquad
 \nabla_{\dot\gamma}(df)=\tfrac12\Delta_h f\,h(\dot\gamma,\cdot),
 \qquad
 \frac{d(\Delta_h f)}{ds}
   =-2\Bigl(\Lambda+\frac{\Delta_h f}{f}\Bigr)df(\dot\gamma).
\]
In a parallel frame this is smooth where $f>0$. Initial data
$(f,0,0)$ give the constant solution by uniqueness. Consequently
$f$ is constant across the local disks, so the box is locally
Nariai by \eqref{eq:einstein-warp}, and $p\in\mathcal P$.
This shows $\mathcal P$ equals the whole connected cover. No
analyticity or compactness of the general late cut has been used.

\section{Global identification}

The cyclic cover is locally symmetric since its metric is locally
a product of constant-curvature factors. Its universal cover
$\widetilde M$, also the universal cover of the original $M$,
is simply connected and null complete. Put
$N=\widetilde{\dS_2(\Lambda)}\times S^2(\Lambda)$, which is
simply connected and geodesically complete. The local product
identifies the tangent curvature model of $\widetilde M$ with that
of $N$ at chosen points. By Lafuente L\'opez
\cite[Theorem~1]{Laf}, a curvature-preserving tangent isometry
between connected locally symmetric null-complete Lorentz
manifolds extends to a local isometry defined on the entire
source if the source is simply connected. Apply this in both
directions to the inverse tangent isometries. Their compositions
have the identity value and differential at their base points;
a local isometry with this initial data is the identity, by
preservation of the geodesic exponential map and continuation
along connected paths. The global maps are therefore inverses.
Equivalently, \cite[Corollary~2]{Laf} first gives geodesic
completeness from null completeness after establishing local
symmetry, and one can then use the complete simply connected
Cartan--Ambrose--Hicks theorem. This proves Theorem~\ref{thm:main}.

The topological and foliation rigidity results of Ling, Rossdeutscher,
Simon and Steinbauer \cite{LRSS} give a complementary comparison
under 2-convexity and additional symmetry hypotheses in parts of their
analysis. The present conclusion concerns the Lorentzian universal-cover
metric for the stated semidefinite vacuum stock.

\begin{thebibliography}{9}
\bibitem{LRSS}
E.~Ling, C.~Rossdeutscher, W.~Simon and R.~Steinbauer,
Rigidity aspects of a cosmological singularity theorem,
\href{https://arxiv.org/abs/2508.00524}{arXiv:2508.00524}.

\bibitem{GLcos}
G.~J.~Galloway and E.~Ling,
A simplified proof of a cosmological singularity theorem,
\emph{Gen.\ Relativ.\ Gravit.} \textbf{58}, 66 (2026).
\href{https://doi.org/10.1007/s10714-026-03562-x}{doi:10.1007/s10714-026-03562-x}.

\bibitem{GLpen}
G.~Galloway and E.~Ling,
Rigidity aspects of Penrose's singularity theorem,
\emph{Commun.\ Math.\ Phys.} \textbf{406}, 25 (2025).
\href{https://doi.org/10.1007/s00220-024-05210-4}{doi:10.1007/s00220-024-05210-4}.

\bibitem{GM}
G.~J.~Galloway and A.~Mendes,
Rigidity of marginally outer trapped 2-spheres,
\emph{Commun.\ Anal.\ Geom.} \textbf{26}, 63--83 (2018).
\href{https://arxiv.org/abs/1506.00611}{arXiv:1506.00611}.

\bibitem{NP}
R.~\v{S}varc, A.~Pravdov\'a and D.~Mi\v{s}kovsk\'y,
Newman--Penrose formalism in quadratic gravity,
\emph{Phys.\ Rev.\ D} \textbf{107}, 024036 (2023).
\href{https://doi.org/10.1103/PhysRevD.107.024036}{doi:10.1103/PhysRevD.107.024036}.

\bibitem{Laf}
J.~Lafuente L\'opez,
A geodesic completeness theorem for locally symmetric Lorentz manifolds,
\emph{Rev.\ Mat.\ Univ.\ Complut.\ Madrid} \textbf{1}, 101--110 (1988).
\href{https://doi.org/10.5209/rev_REMA.1988.v1.n1.18187}{doi:10.5209/rev\_REMA.1988.v1.n1.18187}.

\end{thebibliography}
\end{document}
